% !TeX root = ../papers/nc-p3.tex
\section{中央二次元の中心の対と曲線列の内部での中央支持}
\label{sec:branch-binary-center-pairing}
\begin{theorem}\label{thm:branch-binary-center-pairing}
曲面射影のない標準正則窓で、中央支持2の順中心 $p$（位置0）には
位置3の中央支持2の逆中心 $z$ が一意に対応し、逆も成り立つ。
$p$ の順・中央支持を $S\subset V_1,M\subset V_2$ とすると、$z$ の中央・逆支持も
$S,M$ である。$p$ の後続直線と $z$ の前点直線の一般点は中央支持3である。
逆に分岐中心間の有限曲線列で中央支持が一般に3の動く成分は、
このような対の最初または最後の直線に限る。その他の動く一般点は全て中央支持4である。
反対階数の制限はない。
\end{theorem}
\begin{proof}[二次元の中心からの逆中心の構成]
$w(p)$ の三支持は順に $S,M,C$、次元は $2,2,3$ である。
最後の縮約像 $U\subset R_1\cap(S\otimes M)$ は三次元で、
補題\ref{lem:middle-binary-pencil}の $L_2$ により二項対称二乗の平面となる。
位置2の $x\in N=\mathbf P(M^\perp)$ を一般に取る。$R_1$ を末尾で $x$ に評価すると
$U$ が消えるので逆支持は3以下、動く中心がないので3である。従って一意な前点 $v$ がある。
$v$ が動けば逆支持は3以上で、一定なら $N$ を後続に持ち逆支持3である。
隣接六空間での直交により $x$ の順支持は3以下、再び一般点でちょうど3である。
$x$ の一意な次点を $z$ とする。$z$ が動けば同じ直交により順支持3、
一定なら逆中心なので順支持3である。
$w(z)$ の末尾は $\ker x$ に支持される。その先頭で縮約した三次元像は
$R_1\cap(V_1\otimes\ker x)$ に入る。この評価核は次元3なのでちょうど $U$ に等しい。
従って $w(z)$ の中央・逆支持は $S,M$、特に逆支持2である。
逆中心は有限個なので $z$ は一定、$N$ がその前点直線全体となる。
順逆を入れ替えても全く同じ構成で、$S^\perp$ の射影直線の一意な前点が元の $p$ である。
$w(x,z)=0$ が $p$ の後続直線 $x$ 上で成り立つ（$z$ の中央支持は $S$）。
従って一般点の中央支持は3以下、両端3のため3である。$N$ 上も同じである。
\end{proof}

\begin{proof}[二つの動く反対点]
分岐列の動く一般点 $x$ を位置0とし、その中央支持が3なら唯一の反対消去点を
$y$（位置2）とする。$w(x,y)=0$ であり、両隣の六空間の直交と
$x$ の両端支持3から $y$ の両端支持は3以下である。
$y$ が一定ならその中央支持は動く $x$ の包を消すので2以下となる。中央1は連続純関係を生じ、
中央2では両端は異なる次元で $2,3$ か $3,2$ である。従って $y$ は上の対の中心、
$x$ はその中央支持の消去直線上にあり、一般点の一意な前点または次点が一定となる。

$y$ が動くと仮定する。その両端は一般に3で、$w(x,y)=0$ と中央二項鉛筆の補題から中央も3となる。
$x^-,y^-$ を $x,y$ のそれぞれの前点とする。$w(x)$ の先頭を縮約した三次元関係像は
$R_2\cap(\ker y\otimes V_3)$ であり、その二支持は $\ker y,\ker x^-$ である。
実際包含は三支持から従い、$y$ の順評価の階数3で等号となる。
$w(y^-)$ の先頭も $\ker y$ に入る。その末尾支持は3である。$y^-$ が動けば
逆支持は3以上で、$y$ の順支持3との直交により3以下。定点なら順中心で逆支持3である。
従ってその末尾縮約像が上と同じ三次元関係空間を張る。
ゆえに $y^-$ の順支持は3で、$y^-$ は動く（一定なら順中心となる）。
その中央支持は $\ker x^-$、従って $w(y^-,x^-)=0$ である。
$x,y$ を交換すれば $x^-$ も動き、両者の全三支持が3となる。
この操作を反復すると最初の $x$ の前点を何回取っても動く。
しかし分岐列を逆にたどれば元の順中心という定点に達する。従ってこの場合は起こらない。
\end{proof}

\begin{corollary}\label{cor:rank-nine-nonlinear-corridors}
反対階数9の曲線分岐の各中心の中央支持は2または3である。
どの列でも非直線の連続する成分の一区間内に、射影三空間全体を張る成分は高々一本である。
平面に入る各整成分の次数は3以下である。
\end{corollary}
\begin{proof}
中央支持が3以上の中心では最初の直線の一般中央支持は前定理より4である。
定理\ref{thm:branch-rank-one-residual}より中心の中央支持は $9-6=3$ 以下となる。
連続する非直線成分 $C_i,C_{i+1}$ の一般点の中央支持はともに4である。
同じパラメータの直前・直後の点 $x_{i-1},x_{i+2}$ で縮約した関係は階数1以下である。
実際 $W_{i-1}\subset W_i^\perp$, $W_{i+2}\subset W_{i+1}^\perp$ の間の
対合は直交する $W_i,W_{i+1}$ の一方をそれぞれ割った余次元 $9-4-4=1$ の対合を通る。
この関係は一般に非零である。零なら $x_{i+2}$ から $x_{i-1}$ へ辺ができ、四点を巡回できる。
両成分が非直線なので列の内部にあり、$x_{i+2}$ の一般点での一意な後続をたどると
同じ動く成分を再訪するか、最終直線の一定後続と矛盾する。
従って純関係がある。その一つを固定し、全パラメータで評価すれば、二因子の少なくとも一つは
対応する成分全体を消す。非直線を消せるときはその平面の式に一意である。
連続する純関係の両方に同じ文字は使えないので、非直線成分 $k$ 本の区間の
$k-1$ 辺には異なる $k-1$ 本の平面成分が必要である。

最後に次数4以上の整平面成分を考える。その平面上で順側の四次小行列式は全て
曲線の式で割れる。恒等零なら平面全体が射影像に入り、次数4の場合に限られる。
この場合全小行列式は一つの四次式の定数倍で、一般の評価での
六関係の二次元核は一定である。従って平面上で恒等的に消える非零関係があり、
その平面の式を始文字とする純関係となる。逆側からも同じ平面の式を末文字とする
純関係ができ、連続純関係の禁止に反する。ここは反対階数の制限を使わない。
\end{proof}
